\begin{center}
\begin{tikzpicture}[scale=1]

  

    % Whitebox System (3D effect)
   \node[draw, rectangle, minimum width=2cm, minimum height=1cm] (input) at (-3,0) {$UDG(V,E)$};

    % Whitebox System (3D effect)
    \node[draw, fill=gray!20, minimum width=2.5cm, minimum height=1.5cm] (whitebox) at (0.5, 0) {
        \begin{tikzpicture}[scale=0.75]
            \fill[gray!30] (0,0) rectangle (2.5, 1.5); % Front face
            \fill[gray!20] (0,0) -- (1,0.5) -- (1, 2) -- (0, 1.5) -- cycle; % Left face
            \fill[gray!15] (2.5,0) -- (3.5,0.5) -- (3.5, 2) -- (2.5, 1.5) -- cycle; % Right face
            \fill[gray!15] (1,2) -- (3.5, 2) -- (2.5,1.5)  -- (0, 1.5) -- cycle; % Upper face
            \draw (0,0) rectangle (2.5, 1.5); % Front edge
            \draw (0,0) -- (1,0.5) -- (1, 2) -- (0,1.5); % Left edges
            \draw (2.5,0) -- (3.5,0.5) -- (3.5, 2) -- (2.5,1.5); % Right edges
            \draw (1,0.5) -- (3.5,0.5);
            \draw (1,2) -- (3.5,2);
            
        \end{tikzpicture}
    };

    % Blackbox System (3D effect)
    \node[draw, fill=gray!20, text=white, minimum width=2.5cm, minimum height=1.5cm] (blackbox) at (6, 0) {
        \begin{tikzpicture}[scale=0.75]
            \fill[gray!200] (0,0) rectangle (2.5, 1.5); % Front face
            \fill[gray!200] (0,0) -- (1,0.5) -- (1, 2) -- (0, 1.5) -- cycle; % Left face
            \fill[gray!200] (2.5,0) -- (3.5,0.5) -- (3.5, 2) -- (2.5, 1.5) -- cycle; % Right face
            \fill[gray!200] (1,2) -- (3.5, 2) -- (2.5,1.5)  -- (0, 1.5) -- cycle;
            \draw (0,0) rectangle (2.5, 1.5); % Front edge
            \draw (0,0) -- (1,0.5) -- (1, 2)-- (0,1.5); % Left edges
            \draw (2.5,0) -- (3.5,0.5) -- (3.5, 2)-- (2.5,1.5); % Right edges
            \draw (1,0.5) -- (3.5,0.5);
            \draw (1,2) -- (3.5,2);
        \end{tikzpicture}
    };

    % Output Node
    \node[draw, rectangle, minimum width=2cm, minimum height=1cm] (output_black) at (6, -3) {Result};

    % Draw arrows
    \draw[->] (input) -- (whitebox);
    \draw[->] (whitebox) -- (blackbox) node[midway, below]{diametre = 2};
    \draw[->] (whitebox) -- (blackbox) node[midway, above]{$UDG(V,E)$};
    \draw[->] (blackbox) -- (output_black);
    \node at (0.5,0) {$CDS-UDG$};
    \node at (6,0) {\textcolor{white}{$UDC_1$}};


\end{tikzpicture}
\end{center}

% The CDS-UDG implies that if u can dominate to a set of vertices V, it means that the vertices are UDC covered by a disk centered at u and diameter 2. Hence if G has k size Dominating set, then it means that the set of v vertices can be covered by k disc. Is my explanation for the reduction correct?